\section{Conclusion}
\label{sec:conclusion}

The Steersman is a general-purpose topology programmer for neural
network adapter bridges---but a selective one. Given a pair
specification derived from the co-axial structure of a polytope, it
compiles the corresponding block-diagonal partition with coupling ratios
exceeding 40{,}000:1. Given a specification lacking geometric
coherence, it drives the bridge into connectivity collapse. The
selectivity is the finding: it confirms that successful topology
programming encodes genuine geometric information, not an artifact of
the loss formulation.

Three polytope geometries are confirmed. The octahedron ($n{=}4$)
produces $2{+}2$ blocks at 473{,}622:1---the programme's strongest
signal. The rhombic dodecahedron ($n{=}6$) produces $3{+}3$ blocks at
70{,}404:1. The tesseract ($n{=}8$) produces $4{+}4$ blocks at
41{,}564:1, extending topology programming to four dimensions. A
24-cell experiment ($n{=}24$, D4 root polytope) confirms a
$12{+}12$ eigenvalue split with co/cross $35{,}808$:1 (post-hoc
PC-001 recovery) and Fiedler stabilized at $5.6 \times 10^{-4}$
for 2{,}100 steps (fixed $c_w{=}0.1$).
In every
confirmed case, the Steersman's control laws, loss formulation, and
hyperparameters are identical; only the pair specification changes.

The spectral attractor at Fiedler $\approx 0.09$ is the universal
unprogrammed state. Spectral regularization alone produces this
equilibrium across $n \in \{3, 4, 8, 12\}$ with 10.4\% variation, and
the emanation architecture converges to the same attractor through a
different mechanism. Contrastive loss breaks this equilibrium by
redirecting connectivity from isotropic distribution to specified
axes---a $1{,}130\times$ Fiedler drop at $n{=}6$.

Geometric coherence is required. Wrong-labels (random partition)
produces a clean $3{+}3$ eigenvalue split but co/cross
$\approx 0$---structure without direction. Resonance
(number-theoretic pairs) produces a chain-like eigenvalue pattern with
neither blocks nor direction. These two negative controls, together
with the three positive results and the spectral attractor baseline,
define four cleanly separated regimes that classify all observed bridge
behavior.

Bridge initialization independence confirms that the topology is
determined by the pair specification, not the initial conditions:
three maximally different corpus-coupled initializations converge to
identical block-diagonal structure
(50--52{,}000:1, Frobenius distance ${\sim}10^{-4}$), with the
Steersman actively correcting 75\% wrong initial signs to the
canonical pattern. Yet this topology requires continuous maintenance:
removing the Steersman causes signs to collapse to random within 100
steps and co/cross to decay exponentially toward the isotropic
equilibrium. The block-diagonal configuration is a Steersman-maintained
fixed point---accessible from any initialization, but unstable under
the language modelling loss alone.

Steersman annealing reveals that the optimal operating point is not at
full strength. Linearly reducing the contrastive weight from $0.10$ to
zero produces five regimes: maintenance, interference (co/cross floor
at 7:1), exponential growth (peak 18{,}671:1 at $c_w = 0.017$), gradual
decay, and a cliff (76\% drop) at $c_w = 0$. The topology programmer
works best at approximately 2\% of its nominal weight---full strength
operates in the interference regime where contrastive and LM gradients
conflict. Even at zero weight, the bridge retains 2{,}942:1,
two orders of magnitude above the isotropic equilibrium, confirming
that annealing imprints topology that partially persists beyond the
Steersman's removal.

The practical implication is precise: adapter topology is now a
designable property, not an emergent accident. The pair specification is
the program, the Steersman is the compiler, and the bridge is the
substrate. The compiler accepts real geometry and rejects everything
else---and it operates most efficiently at near-whisper intensity.
